% GENERATED FILE. Edit canonical lessons, then run npm run generate:books.
% canonical-source-hash: fnv1a64:ff6408c77bd0a835
% canonical-lessons: IT-C479-allegria, IT-C479-risata, IT-C479-ridacchiare, IT-C479-entusiasmo, IT-C479-commuoversi

\chapter{Ridere di gusto}
\label{ch:it-ridere-di-gusto}
\hlchaptermodality{479}

\begin{chapteropening}
\textbf{By the end of this chapter:} \emph{I can describe joy, laughter and being moved.}

Complete the last lesson of chapter 479: I can describe joy, laughter and being moved.
\end{chapteropening}

\section[l'allegria]{l'allegria — cheerfulness, joy}

\label{lesson:IT-C479-allegria}



\begin{quote}
Before the new one: say the Italian for envy, then the Italian for to envy.
\end{quote}

\subsection*{You'll want to know: l'allegria}
\textbf{l'allegria} — \textquotedblleft{}cheerfulness, joy\textquotedblright{}.

\begin{quote}
La musica porta allegria alla festa.
\end{quote}

\begin{quote}
\emph{Music brings cheer to the party.}
\end{quote}

Built on \emph{allegro} \textquotedblleft{}cheerful\textquotedblright{}.

\begin{cousinweb}
\emph{allegria} is built on \emph{allegro} \textquotedblleft{}cheerful\textquotedblright{}, usually traced to Vulgar Latin \emph{alecrus}, from Latin \emph{alacer} \textquotedblleft{}lively, quick\textquotedblright{} (Lewis \& Short).
\end{cousinweb}

\begin{grammarlens}[title={What you've built}]
One more word for telling, explaining and describing at B1.
\end{grammarlens}

\subsection*{Your turn}
\begin{itemize}
\raggedright
  \item \emph{Say it:} \emph{l'allegria}
  \item \emph{Say it:} \emph{l'allegria}, once more
  \item \emph{Say it:} a sentence of your own with \emph{l'allegria}
  \item \emph{Recall:} say the Italian for envy, then the Italian for to envy, then say \emph{l'allegria} again
\end{itemize}

\begin{tcolorbox}[breakable,colback=teal!4,colframe=teal!35!black,title={Before you move on}]
What does l'allegria mean? (\textquotedblleft{}Cheerfulness, joy\textquotedblright{}.) Say it once more.
\end{tcolorbox}
\section[la risata]{la risata — laugh, burst of laughter}

\label{lesson:IT-C479-risata}



\begin{quote}
Before the new one: say the Italian for to envy, then the Italian for cheerfulness.
\end{quote}

\subsection*{You'll want to know: la risata}
\textbf{la risata} — \textquotedblleft{}laugh, burst of laughter\textquotedblright{}.

\begin{quote}
Abbiamo fatto una bella risata.
\end{quote}

\begin{quote}
\emph{We had a good laugh.}
\end{quote}

\emph{Fare una risata} — to have a laugh.

\begin{cousinweb}
\emph{risata} is built on \emph{riso} \textquotedblleft{}laughter\textquotedblright{}, which is inherited from Latin \emph{rīsus} (Lewis \& Short), from \emph{rīdēre} \textquotedblleft{}to laugh\textquotedblright{}; \emph{-ata} makes one burst of it.
\end{cousinweb}

\begin{grammarlens}[title={What you've built}]
One more word for telling, explaining and describing at B1.
\end{grammarlens}

\subsection*{Your turn}
\begin{itemize}
\raggedright
  \item \emph{Say it:} \emph{la risata}
  \item \emph{Say it:} \emph{la risata}, once more
  \item \emph{Say it:} a sentence of your own with \emph{la risata}
  \item \emph{Recall:} say the Italian for to envy, then the Italian for cheerfulness, then say \emph{la risata} again
\end{itemize}

\begin{tcolorbox}[breakable,colback=teal!4,colframe=teal!35!black,title={Before you move on}]
What does la risata mean? (\textquotedblleft{}Laugh, burst of laughter\textquotedblright{}.) Say it once more.
\end{tcolorbox}
\section[ridacchiare]{ridacchiare — to giggle, to snigger}

\label{lesson:IT-C479-ridacchiare}



\begin{quote}
Before the new one: say the Italian for cheerfulness, then the Italian for laugh.
\end{quote}

\subsection*{You'll want to know: ridacchiare}
\textbf{ridacchiare} — \textquotedblleft{}to giggle, to snigger\textquotedblright{}.

\begin{quote}
Le ragazze ridacchiavano in fondo alla classe.
\end{quote}

\begin{quote}
\emph{The girls were giggling at the back of the class.}
\end{quote}

Built on \emph{ridere} with \emph{-acchiare}, an ending that makes the action small and repeated.

\begin{cousinweb}
\emph{ridacchiare} is built on \emph{ridere} \textquotedblleft{}to laugh\textquotedblright{}, inherited from Latin \emph{rīdēre} (Lewis \& Short), with the ending \emph{-acchiare} for a small, repeated action.
\end{cousinweb}

\begin{grammarlens}[title={What you've built}]
One more word for telling, explaining and describing at B1.
\end{grammarlens}

\subsection*{Your turn}
\begin{itemize}
\raggedright
  \item \emph{Say it:} \emph{ridacchiare}
  \item \emph{Say it:} \emph{ridacchiare}, once more
  \item \emph{Say it:} a sentence of your own with \emph{ridacchiare}
  \item \emph{Recall:} say the Italian for cheerfulness, then the Italian for laugh, then say \emph{ridacchiare} again
\end{itemize}

\begin{tcolorbox}[breakable,colback=teal!4,colframe=teal!35!black,title={Before you move on}]
What does ridacchiare mean? (\textquotedblleft{}To giggle, to snigger\textquotedblright{}.) Say it once more.
\end{tcolorbox}
\section[l'entusiasmo]{l'entusiasmo — enthusiasm}

\label{lesson:IT-C479-entusiasmo}



\begin{quote}
Before the new one: say the Italian for laugh, then the Italian for to giggle.
\end{quote}

\subsection*{You'll want to know: l'entusiasmo}
\textbf{l'entusiasmo} — \textquotedblleft{}enthusiasm\textquotedblright{}.

\begin{quote}
Ha cominciato il nuovo lavoro con entusiasmo.
\end{quote}

\begin{quote}
\emph{She started the new job with enthusiasm.}
\end{quote}

The adjective is \emph{entusiasta}.

\begin{cousinweb}
\emph{entusiasmo} is a learned borrowing, usually traced through Late Latin \emph{enthusiasmus} to Greek \emph{enthousiasmós} \textquotedblleft{}inspiration by a god\textquotedblright{}, from \emph{éntheos} \textquotedblleft{}possessed by a god\textquotedblright{} (EtymDB, under French \emph{enthousiasme}).
\end{cousinweb}

\begin{grammarlens}[title={What you've built}]
One more word for telling, explaining and describing at B1.
\end{grammarlens}

\subsection*{Your turn}
\begin{itemize}
\raggedright
  \item \emph{Say it:} \emph{l'entusiasmo}
  \item \emph{Say it:} \emph{l'entusiasmo}, once more
  \item \emph{Say it:} a sentence of your own with \emph{l'entusiasmo}
  \item \emph{Recall:} say the Italian for laugh, then the Italian for to giggle, then say \emph{l'entusiasmo} again
\end{itemize}

\begin{tcolorbox}[breakable,colback=teal!4,colframe=teal!35!black,title={Before you move on}]
What does l'entusiasmo mean? (\textquotedblleft{}Enthusiasm\textquotedblright{}.) Say it once more.
\end{tcolorbox}
\section[commuoversi]{commuoversi — to be moved, to be touched}

\label{lesson:IT-C479-commuoversi}



\begin{quote}
Before the new one: say the Italian for to giggle, then the Italian for enthusiasm.
\end{quote}

\subsection*{You'll want to know: commuoversi}
\textbf{commuoversi} — \textquotedblleft{}to be moved, to be touched\textquotedblright{}.

\begin{quote}
Mi commuovo sempre ai matrimoni degli amici.
\end{quote}

\begin{quote}
\emph{I always get emotional at my friends' weddings.}
\end{quote}

The past participle is \emph{commosso}: \emph{ero commosso} — I was moved.

\begin{cousinweb}
\emph{commuoversi} is the reflexive of \emph{commuovere} \textquotedblleft{}to move, to touch\textquotedblright{}, inherited from Latin \emph{commovēre} \textquotedblleft{}to stir up\textquotedblright{} (Lewis \& Short), from \emph{movēre} \textquotedblleft{}to move\textquotedblright{}.
\end{cousinweb}

\begin{grammarlens}[title={What you've built}]
One more word for telling, explaining and describing at B1.
\end{grammarlens}

\subsection*{Your turn}
\begin{itemize}
\raggedright
  \item \emph{Say it:} \emph{commuoversi}
  \item \emph{Say it:} \emph{commuoversi}, once more
  \item \emph{Say it:} a sentence of your own with \emph{commuoversi}
  \item \emph{Recall:} say the Italian for to giggle, then the Italian for enthusiasm, then say \emph{commuoversi} again
\end{itemize}

\begin{tcolorbox}[breakable,colback=teal!4,colframe=teal!35!black,title={Before you move on}]
What does commuoversi mean? (\textquotedblleft{}To be moved, to be touched\textquotedblright{}.) Say it once more.
\end{tcolorbox}
